\documentclass[14pt,solutions]{ca-sheet}
% Generated by taskgen from sheet "test-legacy" (master seed 4242).
% Do not edit: regenerate with `taskgen regen <manifest.json>`.
\begin{document}
\sheettitle{Тест. Решения}


\clearpage

\tableofcontents

\clearpage
\section{Варианты группы A 1}
% variant A 1/1
\subsection{Вариант 1}
\begin{enumerate}
  \item Решить задачу дискретного логарифмирования методом Шенкса: \[29^x \equiv 26\ (\textrm{mod}\ 83).\] Обосновать выбор длины списков.
    \solution{Порядок $g$ делит $p - 1 = 82$; при длине списков $m = \lceil\sqrt{82}\rceil = 10$ любой показатель $x < m^2$ записывается как $x = im + j$, $0 \le i, j < m$ (а $m^2 \ge p - 1$). Малые шаги $g^j$:\[\begin{array}{c|rrrrrrrrrr}j & 0 & 1 & 2 & 3 & 4 & 5 & 6 & 7 & 8 & 9 \\ \hline g^j & 1 & 29 & 11 & 70 & 38 & 23 & 3 & 4 & 33 & 44\end{array}\]$g^{-1} \equiv 63$, $g^{-m} \equiv 63^{10} \equiv 75$. Большие шаги $h \cdot g^{-im}$ до совпадения:\[\begin{array}{c|rrr}i & 0 & 1 & 2 \\ \hline hg^{-im} & 26 & 41 & 4\end{array}\]Совпадение $h g^{-2 \cdot 10} = g^{7}$, откуда $x = 2 \cdot 10 + 7 = 27$. $\ord(29) = 41$: $29^{41} \equiv 1$, $29^{1} \equiv 29$.\[x = 27\ (\textrm{mod}\ 41).\]}
  \item Вычислить степень $19P$ точки $P = (52, 28)$ эллиптической кривой \[y^2 \equiv x^3 + 88x + 33\ (\textrm{mod}\ 97)\] с помощью тернарного разложения показателя $19$.
    \solution{Тернарное разложение (цифры $-1, 0, 1$, без соседних ненулевых): $19 = 1010\bar 1$ ($\bar 1 = -1$), т.е. $19 = 2^{4} + 2^{2} - 2^{0}$; для цифры $-1$ прибавляем $-P = (52, 69)$. Начиная с $R = P$, для каждой следующей цифры: $R \gets 2R$, затем при ненулевой цифре $R \gets R \pm P$. Формулы: $\lambda = \frac{3x_1^2 + a}{2y_1}$ (удвоение), $\lambda = \frac{y_2 - y_1}{x_2 - x_1}$ (сложение), $x_3 = \lambda^2 - x_1 - x_2$, $y_3 = \lambda(x_1 - x_3) - y_1$; $a = 88$, всё по модулю $97$.\[\begin{array}{l|r|r|c}\text{операция} & kP & \lambda & R \\ \hline 2R & 2P & 91 & (29, 28) \\ 2R & 4P & 83 & (41, 43) \\ R + P & 5P & 78 & (74, 2) \\ 2R & 10P & 55 & (64, 63) \\ 2R & 20P & 12 & (16, 28) \\ R - P & 19P & 20 & (41, 54) \\ \end{array}\]\[19P = (41, 54).\]}
\end{enumerate}
\clearpage
% variant A 1/2
\subsection{Вариант 2}
\begin{enumerate}
  \item Возвести целое число в степень по модулю с помощью бинарного алгоритма: \[115^{33}\ (\textrm{mod}\ 336).\]
    \solution{$33 = 100001_2$. Биты показателя обрабатываются слева направо: $r \gets r^2$, при бите $1$ ещё $r \gets r \cdot a$ (в столбце $e$ --- текущий показатель, $r = a^e$).\[\begin{array}{c|r|r|r}\text{бит} & r^2 & r & e \\ \hline 1 & 1 & 115 & 1 \\ 0 & 121 & 121 & 2 \\ 0 & 193 & 193 & 4 \\ 0 & 289 & 289 & 8 \\ 0 & 193 & 193 & 16 \\ 1 & 289 & 307 & 33 \\ \end{array}\]\[115^{33} \equiv 307\ (\textrm{mod}\ 336).\]}
  \item Вычислить степень $27P$ точки $P = (9, 54)$ эллиптической кривой \[y^2 \equiv x^3 + 44x + 11\ (\textrm{mod}\ 89)\] с помощью бинарного разложения показателя $27$.
    \solution{$27 = 11011_2$. Начиная с $R = P$, для каждой следующей цифры: $R \gets 2R$, затем при ненулевой цифре $R \gets R \pm P$. Формулы: $\lambda = \frac{3x_1^2 + a}{2y_1}$ (удвоение), $\lambda = \frac{y_2 - y_1}{x_2 - x_1}$ (сложение), $x_3 = \lambda^2 - x_1 - x_2$, $y_3 = \lambda(x_1 - x_3) - y_1$; $a = 44$, всё по модулю $89$.\[\begin{array}{l|r|r|c}\text{операция} & kP & \lambda & R \\ \hline 2R & 2P & 76 & (62, 12) \\ R + P & 3P & 16 & (7, 67) \\ 2R & 6P & 26 & (39, 80) \\ 2R & 12P & 16 & (0, 10) \\ R + P & 13P & 84 & (16, 70) \\ 2R & 26P & 77 & (23, 14) \\ R + P & 27P & 48 & (47, 80) \\ \end{array}\]\[27P = (47, 80).\]}
\end{enumerate}
\clearpage
\section{Варианты группы B}
% variant B/1
\subsection{Вариант 1 --- Иванов}
\begin{enumerate}
  \item Решить задачу дискретного логарифмирования методом Шенкса: \[25^x \equiv 56\ (\textrm{mod}\ 113).\] Обосновать выбор длины списков.
    \solution{Порядок $g$ делит $p - 1 = 112$; при длине списков $m = \lceil\sqrt{112}\rceil = 11$ любой показатель $x < m^2$ записывается как $x = im + j$, $0 \le i, j < m$ (а $m^2 \ge p - 1$). Малые шаги $g^j$:\[\begin{array}{c|rrrrrrrrrrr}j & 0 & 1 & 2 & 3 & 4 & 5 & 6 & 7 & 8 & 9 & 10 \\ \hline g^j & 1 & 25 & 60 & 31 & 97 & 52 & 57 & 69 & 30 & 72 & 105\end{array}\]$g^{-1} \equiv 104$, $g^{-m} \equiv 104^{11} \equiv 100$. Большие шаги $h \cdot g^{-im}$ до совпадения:\[\begin{array}{c|rrrr}i & 0 & 1 & 2 & 3 \\ \hline hg^{-im} & 56 & 63 & 85 & 25\end{array}\]Совпадение $h g^{-3 \cdot 11} = g^{1}$, откуда $x = 3 \cdot 11 + 1 = 34$. $\ord(25) = 56$: $25^{56} \equiv 1$, $25^{28} \equiv 112$, $25^{8} \equiv 30$.\[x = 34\ (\textrm{mod}\ 56).\]}
  \item Вычислить степень $11P$ точки $P = (68, 14)$ эллиптической кривой \[y^2 \equiv x^3 + 23x + 8\ (\textrm{mod}\ 71)\] с помощью бинарного разложения показателя $11$.
    \solution{$11 = 1011_2$. Начиная с $R = P$, для каждой следующей цифры: $R \gets 2R$, затем при ненулевой цифре $R \gets R \pm P$. Формулы: $\lambda = \frac{3x_1^2 + a}{2y_1}$ (удвоение), $\lambda = \frac{y_2 - y_1}{x_2 - x_1}$ (сложение), $x_3 = \lambda^2 - x_1 - x_2$, $y_3 = \lambda(x_1 - x_3) - y_1$; $a = 23$, всё по модулю $71$.\[\begin{array}{l|r|r|c}\text{операция} & kP & \lambda & R \\ \hline 2R & 2P & 17 & (11, 32) \\ 2R & 4P & 26 & (15, 6) \\ R + P & 5P & 39 & (18, 19) \\ 2R & 10P & 43 & (38, 44) \\ R + P & 11P & 70 & (37, 26) \\ \end{array}\]\[11P = (37, 26).\]}
\end{enumerate}
\clearpage
% variant B/2
\subsection{Вариант 2 --- Петров}
\begin{enumerate}
  \item Возвести целое число в степень по модулю с помощью бинарного алгоритма: \[911^{22}\ (\textrm{mod}\ 424).\]
    \solution{$a = 911 \equiv 63\ (\textrm{mod}\ 424)$. $22 = 10110_2$. Биты показателя обрабатываются слева направо: $r \gets r^2$, при бите $1$ ещё $r \gets r \cdot a$ (в столбце $e$ --- текущий показатель, $r = a^e$).\[\begin{array}{c|r|r|r}\text{бит} & r^2 & r & e \\ \hline 1 & 1 & 63 & 1 \\ 0 & 153 & 153 & 2 \\ 1 & 89 & 95 & 5 \\ 1 & 121 & 415 & 11 \\ 0 & 81 & 81 & 22 \\ \end{array}\]\[911^{22} \equiv 81\ (\textrm{mod}\ 424).\]}
  \item Вычислить степень $47P$ точки $P = (64, 76)$ эллиптической кривой \[y^2 \equiv x^3 + 42x + 31\ (\textrm{mod}\ 97)\] с помощью тернарного разложения показателя $47$.
    \solution{Тернарное разложение (цифры $-1, 0, 1$, без соседних ненулевых): $47 = 10\bar 1000\bar 1$ ($\bar 1 = -1$), т.е. $47 = 2^{6} - 2^{4} - 2^{0}$; для цифры $-1$ прибавляем $-P = (64, 21)$. Начиная с $R = P$, для каждой следующей цифры: $R \gets 2R$, затем при ненулевой цифре $R \gets R \pm P$. Формулы: $\lambda = \frac{3x_1^2 + a}{2y_1}$ (удвоение), $\lambda = \frac{y_2 - y_1}{x_2 - x_1}$ (сложение), $x_3 = \lambda^2 - x_1 - x_2$, $y_3 = \lambda(x_1 - x_3) - y_1$; $a = 42$, всё по модулю $97$.\[\begin{array}{l|r|r|c}\text{операция} & kP & \lambda & R \\ \hline 2R & 2P & 39 & (35, 85) \\ 2R & 4P & 27 & (77, 42) \\ R - P & 3P & 24 & (47, 96) \\ 2R & 6P & 12 & (50, 62) \\ 2R & 12P & 53 & (90, 49) \\ 2R & 24P & 92 & (39, 84) \\ 2R & 48P & 43 & (25, 33) \\ R - P & 47P & 37 & (19, 92) \\ \end{array}\]\[47P = (19, 92).\]}
\end{enumerate}
\clearpage
\end{document}
